\section*{Additional measurement and calibration details}
\label{sec:measurement-notes}
\paragraph{Earlier control-resource comparison (October 4).}
The earlier CPU/PSS scope is retained in Supplementary Table S1; it differs from the role-classified October~27 Fig.~15 (Tables~S31--S33). All nine earlier within-block DDL=1000 versus 50-ms CPU-contrast intervals include zero. Shared deployment history can affect PSS because KubeEdge controllers persist across blocks while Tide and Ray restart. Resource subtotals use matched offered load; completion fractions at 80 requests/s are reported alongside them.

\paragraph{Earlier process scope and tracing protocol.}
The October~4 broader non-application totals include Tide's scheduler/communication and HTTP adapter, dedicated container shim, and full sampled host containerd; KubeEdge CloudCore/EdgeCore and the two kind containers' internal Kubernetes/runtime processes; or Ray core, launcher, agents, Serve proxy, and replica wrapper. Application worker processes are sampled separately and excluded for all systems. Tide is charged the entire sampled shared containerd process, with no subtraction for other host activity. Other host services and unsampled outer kind shims are excluded. The kind footprint belongs to the tested KubeEdge deployment. These older totals are not pure control-plane measurements. CPU/PSS use arrival-window time weighting. Component values are summed within each block before computing the three-block mean, SD, and exploratory 95\% $t$ interval; paired load/DDL contrasts likewise use block totals. KubeEdge controllers persist across blocks while Tide/Ray restart. The corrected Fig.~15 instead separates management, runtime, serving, deployment auxiliaries and known pre-existing background processes, as detailed in Tables~S31--S33.

With the same worker/CPU budget, 1000-ms DDL, and serial arrivals, 30 randomized paired blocks compare empty and pre-created executor pools. Sessions contain a first request and nine reuse requests (600 total, all successful). Monotonic marks partition admission, resource waiting, executor acquisition, binding/preparation, execution-to-exit, and HTTP response; nonoverlapping stages sum to handler time. Executor identities link nested creation/task-start/INIT spans, and the unchanged worker times kernel service separately. Cached images exclude image-pull cost.

\paragraph{Nested initialization spans.}
The main-text Executor lifecycle figure partitions the directly traced request path. Separate nested marks report on-demand container/task setup of 91.96~ms, task-start invocation of 9.95~ms, and creation-through-INIT of 104.81~ms. First acquisition takes 104.97~ms on demand versus 0.002~ms with a pre-created Executor. First-use binding/preparation is 38.94~ms versus 9.27~ms on reuse. These nested spans overlap; the main figure uses a separate nonoverlapping partition. Images are cached, so image-pull cost is excluded.

\paragraph{Instrumentation and microbenchmarks.}
Ten randomized warm trace-on/off pairs contain 50 measured requests per session after one excluded first request (1,020 requests total). Tracing adds a paired 0.016~ms (95\% interval $-0.25$--0.28) against 30.95~ms without tracing. With 64-byte padding, two-process PAIX/protobuf round trips average 70.67~$\mu$s and 13,031/s, including a minimal SHA256 worker. Candidate filter/sort takes 1.43~$\mu$s for four candidates and 101.71~$\mu$s for 256. The PAIX and candidate-selection microbenchmarks characterize these individual operations.

\paragraph{Calibration and sensitivity protocol.}
The main-text calibration uses 68 candidate weights, three calibration and three disjoint test seeds per workload condition: 4896 runs over 36 traces reused across network profiles. All 4,882,536 decisions were independently replayed. Of 48 predeclared renormalized one-coordinate $\pm30\%$ preference perturbations, 46 preserve the selected weights and two select $(0.1,0.9,0)$. The frozen choice has test objective 0.3124 versus a retrospective test-grid best of 0.3189. Reported test results use the frozen calibration choice; the test-grid best is a retrospective reference. Restricting the calibration grid to $w_u\ge0.1$ is a re-selection on the same calibration traces, not a new simulation; it selects $(0.1,0.9,0)$ for all six profiles and changes held-out goodput from 87.13\% to 86.85\% (test objective 0.3124 to 0.3096).

\paragraph{Historical testbeds.}
The original single-factor plots used an Intel Core i7-8565U @ 1.80~GHz (four physical/eight logical cores). One physical core served system/user processes; three served the experiment. Each request contained ten YOLOX tasks requiring two CPU cores, offered at 1/2/4/8/16 requests/s for 20~s.

The heterogeneous testbed contains three ARM Raspberry Pi~4B things, two x86 UP Xtreme edges, an x86 desktop, and a Dell PowerEdge R750 cloud server. Local RTT is within 1~ms at 1~Gbps; cloud emulation adds 50-ms base delay and 5-ms jitter. Workloads combine YOLOX and FaceNet.


\paragraph{Controlled-placement run accounting.}
The primary image-placement figure averages three fully monitored paired Tide and round-robin runs for each arrival pattern. A control-script-interrupted gradual pair is retained separately, together with a further pair under frozen settings. Including the interrupted pair gives a 1.44$\times$ gradual-arrival goodput ratio, compared with the primary 1.45$\times$. The experiment records retain manifests, warmup, driver settings, network-shaping and routing commands, and per-run configurations. A separate all-remote CPU-proxy diagnostic has its own records and workload.

\paragraph{Native score sensitivity records.}
The selected-weight NJ-only 1000/s burst mean includes the first-run gRPC failure with 16/2000 on-time returns. The other repetitions reach 80.0--80.1\%; a separate same-seed diagnostic reaches 1602/2000. The diagnostic is recorded separately from the five-run mean. Native CPU/RSS records include retained routine histories, in addition to active control state.

\paragraph{Local-capability and endpoint-outage protocols.}
The native C4 check uses the runtime version preceding the no-destination terminal-response fix. Across its three outage cycles, three remote-targeted outage requests produce no return within 1.1~s; six critical remote requests before/after outages and three local requests during outages return matching results. A request with insufficient origin CPU is rejected. These observations measure CPU eligibility, local execution, and terminal-response behavior in that version.

A separate connected--disconnected--recovered HTTP-adapter experiment uses two serial worker kernels to examine local-route availability. Disabling the local route also reduces usable capacity. The response gate discards remote returns while the remote kernel can continue running; this protocol measures routing/completion behavior at the adapter rather than native Tide recovery or migration. Model objectives $R/U$ retain the definitions and measurement status stated in the main text.


\paragraph{Controlled-placement worker and network settings.}
The NJ-k1/NJ-k2 emulation uses two 16-logical-CPU Xeon Platinum 8163 hosts. The two prestarted, node-pinned YOLOX-Nano Pods run on Kubernetes~1.34.2, with requests of two CPUs/512~MiB and limits of two CPUs/1~GiB; each executes one inference at a time. Only the remote path is shaped. Each direction has a 1-Gbit/s cap, 32-ms delay with 5-ms variation, and 0.1\% loss. Eight-stream TCP checks deliver 457--866~Mb/s without jitter/loss and 144--168~Mb/s with them; the configured cap is not the achieved throughput. Retained artifacts include shaping, routing, forwarding, warmup, manifests, and per-run settings. Four JPEG resolutions of one archived image supply the 1,000 requests per run; gradual/burst arrival windows are 270/240~s and the per-request deadline is 135~ms.

\paragraph{Request-binding ablation settings.}
NJ-k1 (16-vCPU Xeon Platinum 8163 VM) assigns twelve application vCPUs (4--15) and four driver vCPUs (0--3). Each request performs ten YOLOX-Nano inferences (416$\times$416 ONNX, Runtime~1.18.1, CPU provider, two intra-operation threads); decoding/preprocessing are timed, postprocessing excluded. CC uses six workers, versus twelve when disabled. PI pairs are disjoint with CC, shared by two workers without it. Five randomized blocks cover 1/2/4/8/12/16/32 requests/s: 210 runs, 45,000 successful measured requests, 7,560 separate warmups. Windows/deadlines are 20/180~s. Traces check affinity, binding and quotas. Pointwise $t$(4) intervals use five run-level statistics. Table~\ref{tab:fig7-4plus12-full} retains all 42 cells; the main figure preselects 1/2/4/8/16/32. The separate six-application-vCPU replay is unchanged.

\paragraph{Restricted calibration model parameters.}
The two local and one remote FIFO executors each have two CPUs; node capacities are 4/6/8 CPUs and 1024/1536/2048~MiB. Measured CPU, memory-access, and I/O kernel scales combine with stipulated node factors, 64/256/32-MiB request tokens, and actual-service multipliers in $[0.8,1.2]$. Placement uses estimated service times; realized service includes these multipliers. The two network-budget profiles are 2/12/35 and 10/40/100~ms. Runs last 30~s, with the remote response gate closed during $[10,20)$; responses are lost while queued work continues without retry or migration. Nominal rates are 10/30/60 requests/s, with steady Poisson arrivals or alternating 5-s phases at 0.5/1.5 times nominal rate. Online coefficients $(w_u,w_l,w_d)$ weight remaining CPU, estimated latency, and configured device preference; offline preferences $\lambda_1,\ldots,\lambda_4$ weight the calibration proxies. The six predeclared profiles are $\lambda=(.25,.25,.25,.25)$ (equal), $(2/7,2/7,1/7,2/7)$ (normalized $1{:}1{:}0.5{:}1$ reference), $(.55,.15,.10,.20)$ (goodput), $(.15,.55,.10,.20)$ (resilience), $(.15,.15,.55,.15)$ (utilization), and $(.15,.15,.10,.60)$ (latency); all sum to one.

\paragraph{Control-cost workload settings.}
The NJ-k1 KVM VM exposes 16 vCPUs on a Xeon Platinum 8163 and uses cgroup v2. Its Go kernel processes 65536 bytes for 100 iterations on CPUs~0--1 with Go parallelism one. The serial baselines are Ray Serve~2.58.0 and direct-Pod HTTP KubeEdge~1.23.1; Tide executes locally.
